\begin{corollary}\label{cor:entries}
Fix the row order, the outermost prefixes and the B\'ezout weights above.
The entries of the right multiplier in \(LJ_n=DV\) are determined
by four integers. For \(P,Q\) write \(k=h(P)\), and let \(q\) be
the length of the longest initial segment of the \(k\) chosen
crossing arcs that belong to pairwise distinct cycles of \(P\cup Q\).
Let \(a,a'\) be the two block counts in \(J_n(P,Q)\) and set
\(\ell=\lfloor(q+1)/2\rfloor\). Then
\[
 V_{P,Q}=X^{a-\ell}Y^{a'-\ell}\,
 \sum_{\substack{r\in I_{k+1}\\r\le q+1}}
   a_{k+1,r}(XY)(XY)^{\,\ell-\lfloor r/2\rfloor}.
\]
\end{corollary}
\begin{proof}
In the weighted sum of \eqref{eq:through-quotient} the surviving
prefixes have length at most \(q\). Applying the identity-box face
count to the first \(q\) arcs gives \(a,a'\ge\ell\), whether or not
\(q+1\) belongs to \(I_{k+1}\). Factoring out the displayed
monomial gives the result.
\end{proof}
Thus at order \(n\) one can precompute at most \(n(n+1)/2\) univariate
polynomials, one for each pair \(1\le q\le k\le n\).
For a particular pair of matchings, a traversal of its alternating
cycles determines \(q\). If the total number of cycles is less than
\(\lfloor(k+1)/2\rfloor\), the sum is empty. In particular a connected
pair gives a zero entry of \(V\) whenever the row height is at least
three. These statements concern the change of basis.
